\documentclass[a5paper,10pt]{article}

% https://github.com/InDevRus/filippov-solutions

% Нумерация страниц
\pagenumbering{gobble}

% Настройка полей
\usepackage{geometry}
\geometry{left=1cm}
\geometry{right=1cm}
\geometry{top=1cm}
\geometry{bottom=1.5cm}

% Вставка таблицы
\usepackage{array}
\usepackage{tabularx}

% Инструменты выравнивания формул
\usepackage{amsmath}
\usepackage{mathtools}

% Диакритика в математике
\usepackage{amsfonts}

% Вычисления размеров на ходу
\usepackage{calc}

% Удобные записи производных
\usepackage[italic=true]{derivative}

% Настройки сносок
\usepackage{enotez}

% Длинные знаки векторов
\usepackage{anyfontsize}
\usepackage[b]{esvect}

% Вставка картинок
\usepackage{graphicx}

% Вставка красной строки
\usepackage{indentfirst}
\setlength{\parindent}{1em}

% Условия в командах
\usepackage{xifthen}

% Луа-скрипты
\usepackage{luacode}

% Настройка команд отдельно в математическом и текстовом режимах
\usepackage{mathcommand}

% Для повторения LaTeX кода
\usepackage{multido}

% Конфигурация отступов формул
\delimitershortfall-1sp
\usepackage{mleftright}
\mleftright

% Растягивание объектов
\usepackage{relsize}

% Обтекание текстом
\usepackage{wrapfig}
\usepackage{subcaption}
\usepackage{floatflt}

% Поддержка ввода кириллицы
\usepackage[english, russian]{babel}

% Настройка корневой позиции для компилятора
\usepackage{import}
\providecommand*{\ProjectRootPath}{..}

\import{\ProjectRootPath/preambles/macros/}{theorems}

\import{\ProjectRootPath/preambles/fonts}{font_setup}

% Знак градуса
\usepackage{siunitx}

\import{\ProjectRootPath/preambles/macros/}{operators}
\import{\ProjectRootPath/preambles/macros/}{redefinitions}

\import{\ProjectRootPath/preambles/fonts}{letters_setup}
\import{\ProjectRootPath/preambles/fonts/adjustments}{custom_superscripts}

\import{\ProjectRootPath/preambles/custom}{formatting}
\import{\ProjectRootPath/preambles/custom}{frames}
\import{\ProjectRootPath/preambles/custom}{list_setup}

% TODO: перевести команды на современный стиль объявления

\begin{document}

\begin{framed}
    Для угла $\phi = \ang{30}$ имеем $\tg{\beta} =
    \tg\br{\alpha \pm \phi} =
    \tg\br{\alpha \pm \ang{30}} =
    \dfrac {\sqrt{3} \tg\alpha \pm 1} {\sqrt{3} \mp \tg\alpha}$.
\end{framed}

$3\pow{x}{2} + \pow{y}{2} = C$.
$6x + 2yy\` = 0$.
$y\` = \tg\alpha = - \dfrac {3x} {y}$.

Имеем
$\tg \beta 
= \dfrac {\sqrt{3} \tg\alpha \pm 1} {\sqrt{3} \mp \tg\alpha}
= \dfrac {3\sqrt{3}x \pm y} {\sqrt{3}x \mp 3y} $.
Введём обозначения
$\dfrac {3\sqrt{3}x + y} {\sqrt{3}y - 3x} = g\br{x; y\br{x}}$ и 
$\dfrac {3\sqrt{3}x - y} {\sqrt{3}y + 3x} = h\br{x; y\br{x}}$.
Итак, получаются уравнения изогональных траекторий вида
$$z\`\br{x} = g\br{x; z\br{x}}
\text{ и } 
z\`\br{x} = h\br{x; z\br{x}}
\text{, то есть } 
z\` = \dfrac {3\sqrt{3}x + z} {3x - \sqrt{3}z}
\text{ и } 
z\` = \dfrac {z - 3\sqrt{3}x} {\sqrt{3}z + 3x}.$$

\begin{remark}
    Ниже приведены решения соответствующих уравнений 
    в параметрической форме 
    $\tsys{x\br{t} = C f\br{t} \\ z\br{t} = C t f\br{t}.}$

    \centering
    \begin{tabularx}{0.8\textwidth} {
            | >{\hsize=0.55\hsize\centering\arraybackslash} X
            | >{\hsize=1.45\hsize\centering\arraybackslash} X |}
        \hline
            Уравнение
            & Функция $\,f = f\br{t}$\\
        \hline
            \rule{0mm}{3.4em}
            $\elevate[0.6em]{
                z\lprime = \dfrac {3\sqrt{3}x + z} {3x - \sqrt{3}z}
            }$
            & $\dfrac 
            {\exp\br{\dfrac {1} {\sqrt{2}} \arctg\br{\dfrac {\sqrt{3}t - 1} {2\sqrt{2}} }}} 
            {\sqrt{\sqrt{3} \pow{t}{2} - 2t + 3\sqrt{3}}}$
            \\[1em]
        \hline
            \rule{0mm}{3.4em}
            $\elevate[0.6em]{
                z\lprime = \dfrac {-3\sqrt{3}x + z} {3x + \sqrt{3}z}
            }$
            & $\dfrac 
            {\exp\br{-\dfrac {1} {\sqrt{2}} \arctg\br{\dfrac {\sqrt{3}t + 1} {2\sqrt{2}} }}}
            {\sqrt{\sqrt{3} \pow{t}{2} + 2t + 3\sqrt{3}}}$
            \\[1em]
        \hline
    \end{tabularx}
\end{remark}

\end{document}

